\documentclass[12pt]{article}
\usepackage[T1]{fontenc}
\usepackage[utf8]{inputenc}
\usepackage{lmodern}
\usepackage{textcomp}
\usepackage{enumitem}
\usepackage{geometry}
\geometry{margin=1in}
\begin{document}
\begin{center}
Answer Key - Science - K-12 Level Assessment 4 \
\small (Answers are ordered exactly as in the question paper.)
\end{center}

\section*{Multiple Choice}
\begin{enumerate}[label=\arabic*.]
\item \textbf{Answer:} A
\\ \textit{Options:} A) lions preying on zebras; B) monkeys using twigs to get food; C) birds following migratory patterns; D) bears opening coolers at campsites
\\ \textit{Note:} The answer "lions preying on zebras" is incorrect as an example of an inherited trait; rather, it illustrates a behavioral adaptation that lions have developed to hunt effectively. Inherited traits are characteristics passed down genetically, such as physical features or instincts, rather than learned behaviors or interactions in an ecosystem.
\item \textbf{Answer:} D
\\ \textit{Options:} A) They are similar in shape.; B) They rotate in the same direction.; C) They contain the same number of stars.; D) They have similar elements.
\\ \textit{Note:} The Milky Way Galaxy and other galaxies in the universe share common elemental compositions, primarily consisting of hydrogen and helium, as these elements were formed during the Big Bang, leading to similar chemical building blocks across galaxies. This shared composition is a fundamental aspect of the universe's structure and development.
\item \textbf{Answer:} C
\\ \textit{Options:} A) are conductors of electricity; B) make complete electric circuits; C) are not conductors of electricity; D) make the electricity move quickly
\\ \textit{Note:} The correct answer is based on the fact that plastic and rubber are insulators, meaning they do not conduct electricity, which prevents electrical current from escaping the wires and protects users from electric shock. This property is essential for ensuring safety in electrical applications.
\item \textbf{Answer:} A
\\ \textit{Options:} A) The distance is extremely large.; B) The distance is a rough estimate.; C) Only light can travel between stars.; D) Light signals were timed reflecting from the star.
\\ \textit{Note:} The light-year is used as a unit of measure because it represents the vast distances in space, with one light-year equating to about 5.88 trillion miles, making it more practical for expressing the enormous distance of celestial objects like stars. Since 4.3 light-years is a significant distance relative to everyday measurements on Earth, the light-year provides a more comprehensible way to convey that scale.
\item \textbf{Answer:} C
\\ \textit{Options:} A) turning off lights; B) using public transportation; C) throwing away aluminum cans; D) opening windows to cool a house
\\ \textit{Note:} Throwing away aluminum cans wastes a natural resource because aluminum is a finite material that requires significant energy and resources to mine and process; recycling cans instead conserves these resources and reduces environmental impact.
\end{enumerate}

\section*{True / False}
\begin{enumerate}[label=\arabic*.]
\setcounter{enumi}{5}
\item \textbf{Answer:} False
\\ \textit{Options:} A) True; B) False
\\ \textit{Note:} Recycling centers cannot effectively separate aluminum and steel cans solely by size because both types of cans can be similar in dimensions, and they require different methods, such as magnetic separation for steel and non-magnetic processes for aluminum.
\item \textbf{Answer:} True
\\ \textit{Options:} A) True; B) False
\\ \textit{Note:} A plastic spoon is made from synthetic materials, such as polymers derived from petrochemicals, which are manufactured through chemical processes and do not occur naturally in the environment, thus confirming that the statement is true.
\item \textbf{Answer:} True
\\ \textit{Options:} A) True; B) False
\\ \textit{Note:} Fossil fuels contain hydrocarbons, which are organic compounds that store chemical energy that can be released when they undergo combustion, producing heat and light.
\item \textbf{Answer:} True
\\ \textit{Options:} A) True; B) False
\\ \textit{Note:} The statement is true because the life cycle of a plant typically begins with a seed, which germinates into a young plant and eventually develops into a mature plant, illustrating the natural progression of plant growth stages.
\item \textbf{Answer:} True
\\ \textit{Options:} A) True; B) False
\\ \textit{Note:} Power is indeed measured in watts, which quantifies the rate at which energy is used or transferred, while joules measure the total amount of energy. Since watts and joules represent different aspects of energy-rate versus total amount-they are distinct units.
\end{enumerate}

\section*{Long Form Response}
\begin{enumerate}[label=\arabic*.]
\setcounter{enumi}{10}
\item \textbf{Answer:} As Alex watches the sun set behind the neighbor's house, he observes the effects of Earth's rotation and how it affects light and shadow. The Earth spins on its axis, causing the sun to appear to move across the sky, which creates shadows that lengthen as the sun descends. When the sun dips below the horizon, the light from it is blocked by the house, casting a shadow and creating the beautiful colors of sunset in the sky. This phenomenon illustrates the relationship between light, shadow, and the Earth's movement.
\\ \textit{Note:} The answer correctly identifies that the Earth's rotation causes the apparent movement of the sun across the sky, leading to the creation of shadows that lengthen as the sun sets, while also noting how the house blocks sunlight, resulting in shadows and the vibrant colors of sunset. This demonstrates the interconnectedness of light, shadow, and Earth's motion in the context of the observed phenomenon.
\item \textbf{Answer:} The output energy of a machine is always less than the input energy due to several factors, including friction and energy transformation. Friction between moving parts converts some energy into heat, which is unusable for work. Additionally, not all input energy is transformed into useful work; some is lost during the energy conversion process, leading to a decrease in overall efficiency.
\\ \textit{Note:} correct because it identifies friction as a major source of energy loss, converting usable energy into heat, and emphasizes that not all input energy is effectively transformed into work, thereby reducing the machine's overall efficiency.
\end{enumerate}

\vfill
\noindent\textit{End of Answer Key}
\end{document}